\section{Related work}

\subsection{Approaches to very high spatial resolution land cover classification} 

\note{maybe define VHSR in the intro???}

% VHSR land cover classification has historically been tricky due to blah blah blah
% Prior to deep learning, OBIA was the predominant paradigm - mention a few models
% Performance of OBIA at VHSR attributable to integration of spatial information into imagery classification
% Drawbacks, scale parameters, generalization, human intervention

Characterizing Earth surface in a data-driven manner has long been a central topic of research in the remote sensing community~\citep{}.
Early efforts into automated land cover classification algorithms emerged in the 1970s and 1980s, coinciding with the early days of satellite remote sensing~\citep{}.
Many of these efforts utilized parametric models to assign land cover classes to pixels in medium resolution imagery (e.g., Landsat, SPOT) based on their spectral signatures~\citep{}.
Land cover classification methods matured in the 1990s and early 2000s as computational resources became more widely available and non-parametric machine learning models were adopted for these pixel-wise classification tasks~\citep{}.
Despite the increased predictive capacity of these non-parametric models, classification of very high spatial resolution (VHSR) imagery (e.g., <5m GSD) remained a challenging task due the inherent complexity and heterogeneity of land cover at these finer spatial scales~\citep{}.
\citep{} demonstrated that local spectral variability in land cover classes substantially increases as the spatial resolution of imagery increases beyond 5 m, increasing the inherent intra-class variance and making pixel-wise classification of VHSR imagery inherently challenging.
Object-based image analysis (OBIA) methods emerged in the early 2000s in an effort to address these challenges by grouping pixels into objects using a segmentation algorithm, then classifying these objects as a whole rather than classifying individual pixels~\citep{}.
One notable side effect of the OBIA paradigm is that it effectively introduced spatial context into the classification process, both in an indirect manner by grouping pixels into segments using some and optionally in an explicit manner by incorporating shape descriptors into the feature set passed.
The success of OBIA methods in VHSR land cover classification has been largely attributed to this direct consideration of spatial context, which is difficult to achieve with pixel-wise classification methods, even with intensive feature engineering to extract spatial features from the imagery using filters, texture measures.

However, despite the superiority of OBIA methods over their pixel-wise predecessors, these approaches have several notable drawbacks that have motivated research into alternative methods for classification of VHSR imagery.
Perhaps most obvious is the classification model's reliance on the preceding segmentation algorithm producing yielding groups of pixels that meaningfully delineates objects in the source imagery in an accurate and reliable fashion.
Calibrating a segmentation algorithm to accomplish this is an ill-defined process that requires substantial human supervision, even before a classifier can be trained~\citep{}.
Even if the so-called ``scale parameter'' for a given algorithm is chosen appropriately for a given area, the generalization capacity of an OBIA model is an oft-cited bottleneck~\citep{}.
OBIA models calibrated and trained over rural areas tend to fare poorly in urban environments due to the decreased size of objects found in highly-developed locales, while models developed for urban imagery classification similarly struggle in less developed regions~\citep{}.
While multi-scale segmentation algorithms attempt to address these shortcomings, they introduce additional complexity into the classification process that makes developing and scaling these models for operational land cover mapping difficult.

Subsequently, the introduction of deep learning-driven convolutional neural networks (CNNs) for general-purpose image processing tasks in the broader computer vision research community sparked substantial interest from researchers in remote sensing looking to consolidate the multi-stage classification process into an end-to-end modelling framework capable of incorporating both spectral and spatial features into predictions.
CNNs utilize stacked convolutional kernels with learnable parameters to extract spatially-organized feature maps corresponding to the presence of spectral and spatial features in the input.
Orignally designed for image classification, \citet{longFullyConvolutionalNetworks2015} demonstrated that CNNs could generate dense pixel-wise classification maps by upsampling these feature maps to match the resolution of the input, then applying a simple linear classifier to each pixel in the upsampled feature maps.
Despite the relative simplicity of this approach compared to existing methods, these fully convolutional networks (FCNs) achieved state-of-the-art performance on difficult general-purpose semantic segmentation tasks.
The design of FCNs became a starting point for research into deep learning-models designed explicitly for semantic segmentation, and works centered around integrating these models for classification of remotely sensed data soon followed.
\note{NOTABLE EXAMPLE 1}
\note{NOTABLE EXAMPLE 2}
\note{NOTABLE EXAMPLE 3}
Given the capability of semantic segmentation models to 


% Semantic segmentation introduced in computer vision community
% Why it worked well
% Prominence in VHSR classification today
% List a couple of notable works on convolution

\subsection{Training data sampling strategies}

While deep learning semantic segmentation models have emerged as the predominant approach for land cover classification at VHSR, comparatively little work has been published on the selection of samples for training data compared to their pixel-wise counterparts.
Indeed, most research into the applications of these models to VHSR land cover are evaluated against benchmark datasets, such as the ISPRS Potsdam and Vaihingen 2D Semantic Labelling Challenges.
While these datasets are valuable for determining the efficacy of new models and techniques in the context of existing works, they ignore critical considerations required for real-world applications of these models, particularly considering the data-hungry nature of deep learning models and the importance of quality training data.
Training data selection strategies for pixel-wise approaches have been widely studied such to determine best approaches for 

% Mention stratification, classical pixel-wise
% Struggles when applying to patch wise
% Mention a few

\subsection{Remote sensing pre-training and self-supervised learning}

Bring up pre-training, why its important, how transfer learning works
Difficulty in remote sensing pre-training, unsure of domain shift betweem standard ImageNet pre-training and remote sensing imagery
Emergence of self-supervised learning as a new paradigm for pre-training

How this has translated to remote sensing in the form of foundation models
What these foundation models do, why they are built
Problems with foundation models: ViT lacks inductive bias toward imagery
Transformers still not the norm in remote sensing analysis, lack of inductive bias makes very low data fine-tuning difficult
Still, can apply SSL to CNNs using contrastive/self-distillation approaches